\section{Formal Setting}

We assume the Ontology of Continua (OC Core) as the foundational formalism for
describing admissible system states and structural evolution~\cite{oc_core}.
In addition, we adopt the executable-closed ETS.K\_EA.v1.1 specification as an
immutable symbol discipline and constraint set for the enterprise continuum~\cite{ets_kea_v11}.
No new axioms or structural primitives are introduced in this work.

The purpose of this section is not to rederive OC Core, but to fix the minimal
formal objects required to state and evaluate intervention limits.

\subsection{Enterprise continuum and admissibility}

We represent the enterprise continuum as
\[
K_{EA} \equiv (\Omega,\partial\Omega, A, P, \ThetaSys, J, C, k, \tau),
\]
where $\Omega=\Omega(K_{EA})$ denotes the admissible state space and
$\partial\Omega$ its boundary. The remaining components follow OC Core
conventions: axes $A$, potentials $P$, a threshold system $\ThetaSys$, flows $J$,
closure cycles $C$, a continuity measure $k$, and characteristic cycle times
$\tau$.

Thresholds are treated as a \emph{system} (possibly vector-valued) of inequality
constraints on states:
\[
\ThetaSys(\omega) \le 0
\quad \text{means that $\omega$ is admissible under all thresholds.}
\]
Accordingly, we define the admissible region as
\[
\Omega := \{\omega \in \mathcal{S} \mid \ThetaSys(\omega)\le 0\},
\]
for some ambient space $\mathcal{S}$ of candidate states.
The boundary $\partial\Omega$ separates admissible from inadmissible states and
is the locus at which feasibility becomes fragile.

No smoothness, convexity, or metric structure of $\Omega$ is assumed unless
explicitly stated. In particular, admissibility is governed solely by inequality
constraints rather than by local perturbative notions.

\subsection{Continuity and phase labels}

Structural continuity is represented by a scalar observable
\[
k:\Omega \to [0,1],
\]
interpreted as a measure of connectedness, closure, or structural coherence of
the enterprise continuum, in the OC Core/ETS sense.
We do not assume monotonicity, differentiability, or explicit dynamics for $k$;
it is used only as a comparative indicator before and after intervention.

We further introduce a coarse phase classifier
\[
\Phase:\Omega \to \{\OK,\INERTIA,\COLLAPSE,\STOP\},
\]
which assigns each admissible state to a qualitative regime determined by
OC/ETS-consistent criteria (e.g., closure loss, threshold death, or cycle
divergence).
The analysis does not depend on a specific classifier construction, only on the
existence of a phase partition and the terminal nature of $\STOP$.

\begin{remark}
The boundary $\partial\Omega$ does not represent small perturbations or local
instability.
Rather, it marks structural non-smoothness: near $\partial\Omega$, the set of
admissible mappings changes discontinuously because feasibility is enforced by
inequality constraints.
\end{remark}
